%% Allgemeines -----------------------------------

% Fortlaufende Nummerierung - Entweder 1, 2, 3, 4, ... oder mit einem semantischen Schema
\def \decisionID {19}

% Datum, an dem der Record anefertigt wurde
\def \decisionDate {26.11.2026}

% Wer hat an dem Record mitgewirkt?
\def \decisionParticipants {Quentin Becker, Niklas Käsehagen, OneTutor (Customer Proxy)}

% Der Status in dem sich die Anforderung in diesem Record aktuell befindet.
% Möglichkeiten: Vorgeschlagen, Akzeptiert, Abgelehnt, Abgelöst
\def \decisionStatus {Akzeptiert}

% ID des Records, der diesen Record ablöst, falls zutreffend
\def \decisionSuperseded {}

% ID des Records, den dieser Record ablöst, falls zutreffend
\def \decisionSuperseding {11}

%% Anforderung -----------------------------------

% Titel der Anforderung
\def \requirementTitle {Identische Folien sollen nicht gespeichert werden}

% Kontext der Anforderung
\def \requirementContext {Während der Entwickung der Folien kann es vorkommen, dass zwei Tutor*innen die selbe Änderung vornehmen oder man versucht eine bereits vorhandene Version der Folien erneut hochzuladen. Damit das Repository übersichtlich bleibt sollen komplett identische Versionen nicht gespeichert werden.}

% Konsequenz der Anforderung
\def \requirementConsequences {Ist die XML-Struktur zweier Folien identisch, so wird die, die neu hochgeladen werden soll nicht hochgeladen. Für diese Prüfung ist es nicht nötig das Ähnlichkeitsmaß zu verwenden.}